\documentclass[11pt,letterpaper]{article}
\usepackage[T1]{fontenc}
\usepackage{lmodern,amsmath,amssymb,booktabs,tabularx}
\usepackage[margin=0.91in,headheight=13pt]{geometry}
\usepackage{microtype,fancyhdr,xcolor,enumitem}
\usepackage[hypertexnames=false,colorlinks=true,linkcolor=black,urlcolor=blue!50!black]{hyperref}
\hypersetup{pdftitle={Many-body simulation methods: estimating entropy production from particle counts}}
\pagestyle{fancy}\fancyhf{}
\fancyhead[L]{\small Many-body simulation methods}
\fancyhead[R]{\small Periodically driven particles}
\fancyfoot[C]{\thepage}\renewcommand{\headrulewidth}{0.3pt}
\setlength{\parindent}{0pt}\setlength{\parskip}{5pt plus 1pt}
\setlength{\abovedisplayskip}{7pt}\setlength{\belowdisplayskip}{7pt}
\setlist[enumerate]{leftmargin=*,itemsep=2pt,topsep=1pt,parsep=1pt}
\newcommand{\Tr}{\mathcal T}
\newcommand{\ch}{Q_{\mathrm{ch}}}
\newcommand{\ex}{\hat{\mathbf x}}
\newcommand{\dd}{\mathrm d}
\newcommand{\source}[2]{\href{https://github.com/ruichengbao/measuring_irreversibility_counting/blob/e868697ff37a81a72f4ee47956546a0e9cc1604c/#1}{#2}}
\begin{document}
\thispagestyle{plain}
{\LARGE\bfseries Many-body simulation methods\par}
{\large Estimating entropy production from particle counts\par}
\vspace{3pt}
Technical companion to Fig.~3 and Supplemental Fig.~S3 of
\emph{Estimating Irreversibility via Cross-Correlations and Random Coarse-Graining}.

\section{Measurement and model}
Each frame records the number $n_m(t)$ of particles in every box $m$ of a $K\times K$ grid.
For $N$ particles the normalized counts are $\rho_m(t)=n_m(t)/N$.
Two frames separated by time $\tau$ give the cross-correlation
\begin{equation}
 C^{t_0,\tau}_{nm}=\big\langle\rho_n(t_0)\rho_m(t_0+\tau)\big\rangle .
 \label{eq:C}
\end{equation}
Here $n$ indexes the initial box and $m$ the final box. Repeating the two observations
at the same phase $t_0$ of the drive estimates their ensemble average. Summing over all ordered
box pairs gives $\sum_{n,m}C^{t_0,\tau}_{nm}=1$. Particle identities and trajectories between frames
are unnecessary for this measurement. The simulation evaluates this average by integrating
over the model's known translation statistics. Particle-resolved work provides a separate reference
for the true entropy production (EP).

The model consists of overdamped Brownian particles on a periodic square of side $L=5$.
The unwrapped particle positions $\mathbf x_a$ obey
\begin{equation}
 \dd\mathbf x_a=\mu_a\big[-\nabla_a U+q_a E(t)\ex\big]\,\dd t+\sqrt{2\mu_a}\,\dd\mathbf W_a .
 \label{eq:model}
\end{equation}
Lengths and times use the simulation units, and energies use $k_BT=1$, so that diffusivity
$D_a=\mu_a$ equals mobility. The $\mathbf W_a$ are independent two-dimensional standard Wiener
processes, $\ex$ is the unit vector along the field, and $q_a$ is the particle charge.
The potential $U$ sums two pair energies over $a<b$. At minimum-image distance $r$ and mean
diameter $\sigma_{ab}=(\sigma_a+\sigma_b)/2$, they are
\begin{align}
 V_{\rm WCA}(r)&=\begin{cases}
 4\epsilon[(\sigma_{ab}/r)^{12}-(\sigma_{ab}/r)^6]+\epsilon,&r<2^{1/6}\sigma_{ab},\\
 0,&r\geq2^{1/6}\sigma_{ab},
 \end{cases}\\[-1pt]
 V_{\rm spring}(r)&=\begin{cases}
 \tfrac12 k[(r-r_0)^2-(l_0-r_0)^2],&r<l_0,\\
 0,&r\geq l_0 .
 \end{cases}
\end{align}
The Weeks--Chandler--Andersen repulsion has strength $\epsilon=1$. Every pair has spring
stiffness $k=5$, rest length $r_0=10^{-3}$ and cutoff $l_0=2$. The energy shift makes the
spring potential continuous at the cutoff. Each particle's diameter $\sigma_a$ and charge $q_a$
are drawn independently and uniformly from $\{0.01,0.02,0.03\}$ and $\{0.9,1.0,1.1\}$. The mobility is
$\mu_a=0.01/\sigma_a$.

The field is $E(t)=E[1+\sin(2\pi t/\Tr)]$ with period $\Tr=1$. The simulator takes its
time origin at the field maximum and uses the equivalent cosine protocol. The scan covers
$N=1,2,5,10,20,50,60,80,100$ and $E=0.1,0.2,0.4,0.5,0.8,1.0$. Five realizations at each
size retain the same particle assignments across fields. The
\source{numerics/Fig3_corrected/simulation/jobs_manifest_driven_final.json}{production manifest}
records thermal and disorder seeds and run lengths for the 240 interacting runs.
The 30 one-particle cases use the exact propagator. Equation numbers prefixed by S identify formulas moved from the Supplemental Material.

\newpage
\section{Relative-coordinate simulation and the true EP}
Integrating over the collective displacement removes the sampling noise of the cluster's slow diffusion.
Define the total friction $\Gamma=\sum_a\mu_a^{-1}$, total charge $\ch=\sum_a q_a$,
center of friction $\mathbf R=\Gamma^{-1}\sum_a\mathbf x_a/\mu_a$, and relative positions
$\mathbf y_a=\mathbf x_a-\mathbf R$. They satisfy $\sum_a \mathbf y_a/\mu_a=0$.
Pair forces cancel in the weighted sum of Eq.~\eqref{eq:model}. With
$\mathbf F_a^{\rm int}=-\nabla_a U$ and $\tilde q_a=q_a-\ch/(\mu_a\Gamma)$,
\begin{align}
 \dd\mathbf R&=\frac{\ch}{\Gamma}E(t)\ex\,\dd t+\sqrt{\frac{2}{\Gamma}}\,\dd\mathbf W_R ,
 \tag{S41}\label{eq:cm}\\
 \dd\mathbf y_a&=\mu_a\big[\mathbf F_a^{\rm int}+\tilde q_a E(t)\ex\big]\,\dd t
 +\sqrt{2\mu_a}\,\dd\mathbf W_a-\frac1{\Gamma}\sum_b\sqrt{2/\mu_b}\,\dd\mathbf W_b .
 \tag{S42}\label{eq:rel}
\end{align}
Here $\mathbf W_R=\Gamma^{-1/2}\sum_b\mu_b^{-1/2}\mathbf W_b$ has zero covariance with the
Gaussian noise in Eq.~\eqref{eq:rel}, so the collective increments are independent of the
relative dynamics. The relative mobility matrix is $\mu_a\delta_{ab}-1/\Gamma$, with
$\delta_{ab}$ the Kronecker delta.

\textbf{Numerical update.}
The code integrates Eq.~\eqref{eq:rel} in sweeps with time step $\Delta t=2\times10^{-4}$.
It chooses ascending or descending particle order with equal probability and holds the field
at its midpoint value $E_{\rm mid}=E(t+\Delta t/2)$ during a sweep. For particle $a$, define
$s_a=\sqrt{2\mu_a\Delta t}$ and $\mathbf d_a=\mu_a\Delta t(\mathbf F_a^{\rm int}+\tilde q_a E_{\rm mid}\ex)$.
Propose
\begin{equation}
 \boldsymbol\delta=\frac{\mathbf d_a}{1+|\mathbf d_a|/(2s_a)}+s_a\boldsymbol\eta ,
 \qquad \boldsymbol\eta\sim{\cal N}(\mathbf 0,\mathbf I_2),\tag{S43}
\end{equation}
where $\mathbf I_2$ is the two-dimensional identity matrix. Let $g_{\rm f}$ be this Gaussian
proposal density at $\boldsymbol\delta$ and $g_{\rm r}$ the Gaussian density at
$-\boldsymbol\delta$ using the drift evaluated at the proposed position. For pair-energy change
$\Delta U$, accept the move with probability
\begin{equation}
 \min\!\left\{1,\exp[-\Delta U+\tilde q_a E_{\rm mid}\delta_x]\,g_{\rm r}/g_{\rm f}\right\}.
\end{equation}
The bounded proposal drift controls steps near a steep repulsive core. The acceptance probability
uses the full pair energy and includes external work. After each sweep the code subtracts a common
shift to restore $\sum_a\mathbf y_a/\mu_a=0$. Translation invariance makes this equivalent to
projecting each accepted move. The direct and cached implementations use the same update. Each run
discards 10 initial periods. The full production record spans 160--10,000 periods, with longer
records at small $N$ and weak fields.

\textbf{Reference work.}
In the periodic steady state, the mean changes in potential energy and system entropy vanish
over a cycle. The true EP per period, in units of $k_B$, equals the mean external work
$\langle\sum_a\int q_a E(t)\circ\dd x_{a,x}\rangle$ in units of $k_BT$, with the integral
taken over one cycle (Stratonovich convention). For relative polarization
$\Pi(t)=\sum_a q_a y_{a,x}(t)$, Eq.~\eqref{eq:cm} gives
\begin{equation}
 \sigma_{[0,\Tr]}=\frac{3\Tr E^2\ch^2}{2\Gamma}
   +\left\langle\int_0^\Tr E(t)\,\dd\Pi(t)\right\rangle .
 \label{eq:work}
\end{equation}
The collective term follows from $\int_0^\Tr E(t)^2\dd t=3\Tr E^2/2$. For cycles starting at
a field maximum, integration by parts gives the relative work as
$2E[\Pi(\Tr)-\Pi(0)]-\int_0^\Tr E'(t)\Pi(t)\dd t$. The code retains the boundary term,
compares this expression with direct displacement work and averages over all production periods.
Dividing Eq.~\eqref{eq:work} by $\Tr$ gives the reference rate.

\newpage
\section{Evaluating the two-time count correlations}
Figure~3 uses integer-period lags $\tau=\ell\Tr$ with $\ell=1,\ldots,40$. In the sine convention
the initial phase $t_0$ is a field maximum ($\Tr/4$) or minimum ($3\Tr/4$). The field satisfies
$E(t_0+s)=E(t_0+\tau-s)$ over each window and the periodic steady-state endpoint distributions
agree. The reversed experiment therefore has the same correlation matrix as the forward one.
Reversing the time order exchanges its two box indices. The equal-weight average of the two
extremum matrices also gives a valid bound by joint convexity of relative entropy. General
driving windows require forward and reversed observations.

At any fixed phase the periodic center $\mathbf R\bmod L$ is uniform on the box
$\Omega=[0,L)^2$, independently of the relative configuration. Over an integer-period lag
its unwrapped increment $\boldsymbol\Delta$ is Gaussian with mean
$\bar\Delta\ex=(\ch E\tau/\Gamma)\ex$ and variance $s^2=2\tau/\Gamma$ per axis.
Write $G_\tau$ for this two-dimensional density, $\mathbf y(t)$ for the collection
of all relative coordinates, and
$\rho_m(\mathbf R,\mathbf y)=N^{-1}\sum_a\mathbf 1_m(\mathbf R+\mathbf y_a)$,
where $\mathbf 1_m$ is the periodic box indicator. The conditional average of the same two
counts in Eq.~\eqref{eq:C} is
\begin{equation}
 C^{t_0,\tau}_{nm} =
 \left\langle \frac{1}{L^2}\int_\Omega\dd^2 R
 \int_{\mathbb R^2}\dd^2\Delta\,G_\tau(\boldsymbol\Delta)
 \rho_n(\mathbf R,\mathbf y(t_0))\,
 \rho_m(\mathbf R+\boldsymbol\Delta,\mathbf y(t_0+\tau))\right\rangle_{\mathbf y}.
 \tag{S47}\label{eq:conditional}
\end{equation}
Only the average over relative configurations is sampled. The analysis uses the last 160 production
periods, recording at both field extrema (320 frames), with all time origins that have an endpoint
in the record.

\textbf{Box probabilities.}
Translation symmetry reduces the ordered matrix to the distribution of the box displacement
$d=m-n$ (a two-dimensional index modulo $K$ in each direction), with
$P_d=\sum_n C^{t_0,\tau}_{n,n+d}=K^2 C^{t_0,\tau}_{0,d}$. A box has width $h=L/K$.
Let $g_L(z;s)$ be the one-dimensional Gaussian density of standard deviation $s$, wrapped
to period $L$. Integrating over the initial position and the final box gives the kernel
\begin{equation}
 \kappa_h(z;s)=\frac1h\int_{-h}^h(h-|v|)g_L(z+v;s)\,\dd v .
\end{equation}
For each ordered pair $(a,b)$, including $a=b$, form the relative displacement
$\mathbf u=\mathbf y_b(t_0+\tau)-\mathbf y_a(t_0)$ modulo $L$. Write $H(\mathbf u)$
for its normalized distribution over pairs and time origins. Equation~\eqref{eq:conditional} becomes
\begin{equation}
 P_d=\int_{\Omega}H(\mathbf u)\,
 \kappa_h(d_xh-u_x-\bar\Delta;s)\,\kappa_h(d_yh-u_y;s)\,\dd^2u.
\end{equation}
The Gaussian and box integrals are analytical. The code approximates $H$ by distributing
each sample among the four neighboring grid points with nonnegative bilinear weights.
It averages over reflection in the transverse coordinate and computes the finest observation
grid ($K=160$) first. Merging its boxes gives the other grids exactly through
$C^{\rm coarse}=A C^{\rm fine}A^{\mathsf T}$, where $A_{mi}=1$ if fine box $i$ lies in
coarse box $m$, and zero otherwise.

The integration grid has $S=128$ points per side for $N\le5$, $256$ for $5<N\le20$, and
$512$ for $N>20$ in the interacting runs. $S$ controls numerical precision. The number of
boxes observed along each side is $K\in\{5,10,20,40,80,160\}$. For $N=1$, $\mathbf y_1=0$ and
the code evaluates the wrapped Gaussian propagator directly and scans the lag without trajectory
sampling. The integer-period displacement then depends only on the constant field component.
For the drive in Sec.~1 even a perfectly resolved unwrapped displacement recovers at most $2/3$ of
the translational EP over the period.

\newpage
\section{Evaluating the bound on independent realizations}
Suppressing the phase and lag indices, write $C^{\mathsf T}$ for the transposed count matrix.
The population counting bound per period is $(\Tr/\tau)D(C\Vert C^{\mathsf T})$, where
$D(C\Vert C^{\mathsf T})=\sum_{n,m}C_{nm}\ln(C_{nm}/C_{mn})$.
Selecting the largest value after fitting a noisy matrix on the same trajectory can
bias the estimate upward. For a real matrix $f$ the variational identity reads
\begin{equation}
 D(C\Vert C^{\mathsf T})=\max_f S(C,f),\qquad S(C,f)=\sum_{n,m}\big[C_{nm}f_{nm}-C_{mn}(e^{f_{nm}}-1)\big].\tag{S48}
\end{equation}
Each term is concave in $f_{nm}$ and attains its maximum at $\ln(C_{nm}/C_{mn})$ for
positive entries. At the maximum the exponential terms cancel because both matrices sum to one.
Every fixed $f$ therefore gives a lower bound linear in the measured correlations.

\textbf{Fitting and scoring.}
For each $N>1$, field $E$ and observation grid $K$, each of the five realizations is evaluated once
with a function learned from the other four. The four training histograms use their own
relative coordinates and the evaluated realization's fixed $\ch$ and $\Gamma$ in the Gaussian
kernels. These two constants describe its particle parameters and require no use of its
thermal trajectory during selection.
\begin{enumerate}
\item For each lag and phase choice (maximum, minimum, or the equal-weight average), hold out
one of the four training realizations. Average the probabilities $P_d$ on the other three,
and fit $f_d$ as the log ratio $\ln(P_d/P_{-d})$ clipped to $[-c,c]$.
Try $c=8,12,16,20,24,28,32,40,50,64,80,100,150$.
\item Score each candidate on the fourth training realization using $S(C,f)/\ell$,
with $f_{nm}=f_{m-n}$. Select the lag, phase and cap that maximize the average score
over all four inner holdouts. Refit $f$ using all four training realizations.
\item Evaluate $S(C^{\rm test},f)/\ell$ on the fifth realization using its own correlation
matrix $C^{\rm test}$. Repeat until each realization has provided the test data.
\end{enumerate}
The population inequality carries over in expectation to an independently estimated linear
score, subject to time-step and quadrature error. An individual finite trajectory can fluctuate
on either side of its population value. The single-particle code uses its deterministic
probabilities for every lag (log-ratio cap 20).

\textbf{Figure values.}
The plotted rate is the mean of the five evaluated per-period scores divided by $\Tr$.
Rate error bars give the sample standard deviation across the five realizations. The quality
factor $\mathcal Q$ divides the mean estimated rate by the mean true rate. Figure~3 uses $K=20$. Its inset compares all six grids at $N=5$, $E=0.1$. Refining $K$ from 20 to 160 raises the
estimated rate there by $2.2\%$. At $N=100$, $E=1$ the same refinement raises $\mathcal Q$
from $0.4156$ to $0.4996$. The grids share the same trajectories and each grid has its own
training-only selection.

\textbf{Numerical checks.}
For $N=50,100$ at $E=0.1,1.0$, decreasing $\Delta t$ from $2\times10^{-4}$ to
$5\times10^{-5}$ changed the true rate by at most $1.32\%$ and the $K=160$ score, evaluated
with the same fitted function and window, by at most $2.10\%$. The differences were within
$1.33$ combined trajectory sampling standard errors. Separate $E=0$ controls gave zero work
for $N=20,100$. Selected tests doubling $S$ changed the count-matrix relative entropy by
less than $0.11\%$. These checks apply to the tested cases.

\textbf{Run and inspect the calculations.}
The \source{README.md}{repository README} gives the commands for regenerating trajectories,
rebuilding tables and plotting. The \source{numerics/Fig3_corrected/analysis/analyze_group.py}{selection code},
\source{numerics/Fig3_corrected/analysis/positive_kernel.py}{box integrals},
\source{numerics/Fig3_corrected/validation/validation_summary.json}{validation results} and
\source{numerics/Fig3_grid_update/plotting_tables/fig3_K20_summary.csv}{Fig.~3 data}
are linked to the simulation version used for this note.
\end{document}
